\section{Выводы}

Выполнив лабораторную работу по курсу \enquote{Дискретный анализ}, я реализовал красно-чёрное дерево и на его основе разработал программу-словарь, поддерживающую операции вставки, удаления, поиска, а также сохранения и загрузки данных в бинарном формате.

В процессе работы были изучены принципы функционирования сбалансированных бинарных деревьев поиска, а именно механизмы поддержания их свойств с помощью перекрашивания узлов и поворотов. Особое внимание было уделено корректной реализации балансировки после вставки и удаления элементов, что является одной из наиболее сложных частей алгоритма.

Также были получены практические навыки работы с динамической памятью, рекурсивными алгоритмами и бинарным вводом-выводом. Реализация сериализации и десериализации структуры данных позволила лучше понять способы представления сложных структур в файлах и обработки ошибок при работе с ними.

В ходе тестирования производительности было установлено, что реализованное дерево уступает по скорости контейнеру \texttt{std::map}, однако демонстрирует ожидаемую асимптотическую сложность $O(\log n)$ для основных операций. Это подтверждает корректность выбранного подхода и реализации.

Данная работа позволила закрепить теоретические знания о структурах данных и алгоритмах на практике, а также получить опыт разработки и тестирования сложных программных решений.
\pagebreak